\section{Цель работы}
\begin{itemize}
    \item Исследование вольт-амперной характеристики (ВАХ)
полупроводникового диода
    \item Исследование работы однополупериодного выпрямителя
    \item Исследование работы мостового выпрямителя
\end{itemize}

\section{Исходные данные}

Диод в соответствии с вариантом - \textbf{MUR460}.

\section{Снятие ВАХ полупроводникового диода}
\textbf{Электрическая схема (рис. \ref{el_sch_1}).}\\
\begin{figure}[H]
    \centering
    \includegraphics[width=1\linewidth]{/home/ns/Desktop/ITMO-courses/Electronics/ElSch_lab1/img/s1.png}
    \caption{Схема}
    \label{el_sch_1}
\end{figure}
\textbf{Осциллограмма (рис. \ref{osc_1}).}\\
\begin{figure}[H]
    \centering
    \includegraphics[width=1\linewidth]{/home/ns/Desktop/ITMO-courses/Electronics/ElSch_lab1/img/o1.png}
    \caption{ВАХ}
    \label{osc_1}
\end{figure}

\textbf{Расчёты}\\
Определим статическое и динамическое сопротивление по двум точкам из ВАХ :\\
$R_{пр}=\frac{U_{пр}}{I_{пр}}=\frac{1.3047584}{7.942252}=0.164$ [Ом]\\
$R_{пр}=\frac{\Delta U_{пр}}{\Delta I_{пр}}=\frac{1.7813295 - 1.3047584}{16.976843 - 7.942252}=0.0527$ [Ом]\\

Пороговое напряжение равно $\approx550$мВ.\\
Определим напряжение пробоя (рис. \ref{osc_11}):
\begin{figure}[H]
    \centering
    \includegraphics[width=1\linewidth]{/home/ns/Desktop/ITMO-courses/Electronics/ElSch_lab1/img/o11.png}
    \caption{Определение пробоя}
    \label{osc_11}
\end{figure}
Получаем $800$ В. Это значение больше максимального обратного ($600$ В по документации).



\section{Исследование однополупериодного выпрямителя}
\textbf{Электрическая схема (рис. \ref{el_sch_2}).}\\
\begin{figure}[H]
    \centering
    \includegraphics[width=1\linewidth]{/home/ns/Desktop/ITMO-courses/Electronics/ElSch_lab1/img/s2.png}
    \caption{Схема}
    \label{el_sch_2}
\end{figure}
\textbf{Осциллограмма (рис. \ref{osc_2}).}\\
\begin{figure}[H]
    \centering
    \includegraphics[width=1\linewidth]{/home/ns/Desktop/ITMO-courses/Electronics/ElSch_lab1/img/o2.png}
    \caption{ВАХ}
    \label{osc_2}
\end{figure}

\textbf{Расчёты}\\
Максимальное мгновенное напряжение $\approx 9.5$В.\\
Среденевыпрямленное напряжение:
$U_{\text{ВЫХ ср}} = \frac{1}{T} \int_0^T U_m \sin(\omega t)  dt = \frac{U_m}{\pi}=\frac{9.5}{3.14}=3.025$B.\\
Максимально обратное напряжение равно $\approx 9.5$В (см рис. \ref{f}).
\begin{figure}[H]
    \centering
    \includegraphics[width=1\linewidth]{/home/ns/Desktop/ITMO-courses/Electronics/ElSch_lab1/img/o22.png}
    \caption{ВАХ}
    \label{f}
\end{figure}


\section{Исследование двухполупериодного выпрямителя}
\textbf{Электрическая схема (рис. \ref{el_sch_3}).}\\
\begin{figure}[H]
    \centering
    \includegraphics[width=1\linewidth]{/home/ns/Desktop/ITMO-courses/Electronics/ElSch_lab1/img/s3.png}
    \caption{Схема}
    \label{el_sch_3}
\end{figure}
\textbf{Осциллограмма (рис. \ref{osc_3}).}\\
\begin{figure}[H]
    \centering
    \includegraphics[width=1\linewidth]{/home/ns/Desktop/ITMO-courses/Electronics/ElSch_lab1/img/o3.png}
    \caption{ВАХ}
    \label{osc_3}
\end{figure}
\textbf{Расчёты}\\
Максимальное мгновенное напряжение $\approx 9.7$В.\\
Cредневыпрямленное значение
напряжения на выходе выпрямителя:
$
U_{\text{Вых ср}} = \frac{2U_m}{\pi} \approx 0,9U_m=8.73
$\\
Максимально обратное напряжение равно $\approx 9.06$В (см рис. \ref{f}).
\begin{figure}[H]
    \centering
    \includegraphics[width=1\linewidth]{/home/ns/Desktop/ITMO-courses/Electronics/ElSch_lab1/img/fs.png}
    \caption{ВАХ}
    \label{fs}
\end{figure}



\section{Исследование двухполупериодного выпрямителя с
емкостным фильтром}
\textbf{Электрическая схема (рис. \ref{el_sch_4}).}\\
\begin{figure}[H]
    \centering
    \includegraphics[width=1\linewidth]{/home/ns/Desktop/ITMO-courses/Electronics/ElSch_lab1/img/s4.png}
    \caption{Схема}
    \label{el_sch_4}
\end{figure}
\textbf{Осциллограмма (рис. \ref{osc_4}).}\\
\begin{figure}[H]
    \centering
    \includegraphics[width=1\linewidth]{/home/ns/Desktop/ITMO-courses/Electronics/ElSch_lab1/img/o4.png}
    \caption{ВАХ}
    \label{osc_4}
\end{figure}

\textbf{Расчёты}\\
Максимальное значение напряжения $0.6$Вб минимальное $-9.36$В.\\
Среднее значение $-4.3316$В.\\
Определим коэффициент пульсаций:
$$
k = \frac{U_{\text{ВЫХ}_{max}} - U_{\text{ВЫХ}_{min}}}{U_{\text{ВЫХ}_{cp}}}\approx2.308
$$
Коэффициент пульсаций получился больше чем у двухполупериодного выпрямителя $(k\approx0,67)$.\\
\textit{Анализ}\\
Увеличение емкости не влияет значиительно на $k$, как и изменение частоты.





\section{Выводы}
В ходе лабораторной работы были экспериментально исследованы характеристики и режимы работы полупроводникового диода MUR460.

\begin{enumerate}
    \item \textbf{Вольт-амперная характеристика диода MUR460}:
    \begin{itemize}
        \item Определено пороговое напряжение диода $\approx 550$ мВ
        \item Рассчитано статическое сопротивление: $R_{пр} = 0.164$ Ом
        \item Рассчитано динамическое сопротивление: $r_{пр} = 0.0527$ Ом
        \item Экспериментально определено напряжение пробоя $\approx 800$ В, что превышает паспортное значение 600 В
    \end{itemize}
    
    \item \textbf{Однополупериодный выпрямитель}:
    \begin{itemize}
        \item Получено средневыпрямленное напряжение $U_{вых ср} = 3.025$ В
        \item Определено максимальное обратное напряжение на диоде $\approx 9.5$ В
        \item Подтверждена низкая эффективность однополупериодной схемы
    \end{itemize}
    
    \item \textbf{Двухполупериодный мостовой выпрямитель}:
    \begin{itemize}
        \item Получено средневыпрямленное напряжение $U_{вых ср} = 8.73$ В
        \item Коэффициент пульсаций составил $k \approx 0.67$
        \item Подтверждено преимущество двухполупериодной схемы по эффективности выпрямления
    \end{itemize}
    
    \item \textbf{Двухполупериодный выпрямитель с емкостным фильтром}:
    \begin{itemize}
        \item Наблюдается значительное увеличение коэффициента пульсаций до $k \approx 2.308$
        \item Изменение емкости конденсатора и частоты не оказывает существенного влияния на коэффициент пульсаций в данной конфигурации
        \item Обнаружена инверсия полярности выходного напряжения
    \end{itemize}
\end{enumerate}

\textbf{Общий вывод}: Экспериментальные результаты в целом соответствуют теоретическим ожиданиям. Наблюдаемое превышение напряжения пробоя может быть связано с условиями проведения эксперимента. Мостовая схема выпрямления подтвердила свои преимущества перед однополупериодной по эффективности.